\documentclass{article}
\usepackage[a4paper, total={6.5in, 10in}]{geometry}
\setlength{\parindent}{0pt}
\usepackage{tcolorbox}
\tcbset{colback=white!94!black, colframe=black!60!white, coltitle = white, boxrule=0.8pt, arc=2mm}
\newtcolorbox{boxs}[2][]{title= #2,#1}
\usepackage{datetime2}
\usepackage[british]{babel}
\usepackage{amsfonts}
\usepackage{amsmath}

\title{\textbf{Real Analysis HW 8}}
\author{Robert O'Connell}
\date{\today}
\begin{document}
\maketitle
\vspace{5mm}
\subsection*{BS - Page 144 Exercise 8}
Prove that if \(f\) and \(g\) are each uniformly continuous on \(\mathbb{R}\), then the composite function
\(f \circ g\) is uniformly continuous on \(\mathbb{R}\)
\\

Since \(f,g\) are two uniformly continuous functions
\(\forall \ \epsilon >0 \ \exists \ \delta_0, \delta_1\) s.t. (with each \(\delta\) independent of \(c\))
\[
|x-c| < \delta_0 \Rightarrow |g(x) - g(c)| < \epsilon 
\]
and
\[
|x-c| < \delta_1 \Rightarrow |f(x) - f(c)| < \epsilon 
\]

For \(g\), picking its \(\epsilon = \delta_1\), we get the chain

\[
|x-c| < \delta_0 \Rightarrow |g(x) - g(c)| < \delta_1 \Rightarrow |f(g(x)) - f(g(c))| < \epsilon 
\]
Thus \(f \circ g\) is uniformly continuous, since the choice of both deltas were independent of \(c\)

\subsection*{BS - Page 167 Exercise 7}
Suppose that \(f:\mathbb{R} \rightarrow \mathbb{R}\) is differentiable at \(c\) and that \(f(c) = 0\). Show
that \(g(x) := |f(x)|\) is differentiable at \(c\) if and only if \(f \ '(c) = 0\)
\\

\((\Rightarrow)\)

Suppose \(g\) differentiable at \(c\),

Suppose FTSOC \(f'(c) \neq 0\), 

\(\Rightarrow f(c)\) is not a relative minimum or maximum

\[
\Rightarrow \forall \ \delta > 0  \ |x-c|< \delta \Rightarrow f(c) > f(x) \text{ and } f(c) < f(x)
\]
Thus in some delta interval around \(c\), \(f\) attains both postive and negative values, since \(f(c) = 0\)

Also we have
\[
g \ ' (c) = \lim_{x\to c}{\frac{g(x) - g(c)}{x-c}} = \lim_{x\to c}{\frac{|f(x)|}{x-c}}
\]

Then choosing sequences which converge to \(c\) from the left and the right, say \(x_n = c - \frac{1}{2n}\) and \(y_n = c + \frac{1}{2n}\), 
we have that \(x_n - c < \delta \Rightarrow f(x_n) < 0\) or \(y_n - c < \delta \Rightarrow f(y_n) < 0\), WLOG assume the first case

\[
x_n - c < \delta \Rightarrow f(x_n) < 0 \Rightarrow |f(x_n)| = - f(x_n)
\]
\[
\Rightarrow g \ ' (c) = \lim_{n\to \infty}{\frac{-f(x_n)}{x_n - c}} = - f ' (c)
\]

Then
\[
y_n - c < \delta \Rightarrow f(y_n) > 0 \Rightarrow |f(y_n)| =  f(y_n)
\]
\[
\Rightarrow g \ ' (c) = \lim_{n\to \infty}{\frac{f(y_n)}{y_n - c}} = f \ ' (c)
\]

Hence \(f \ '(c) = - f \ ' (c)\), which is a contradiction since we assumed \(f \ ' (c) \neq 0\)
\\

\((\Leftarrow)\)

Suppose \(f \ '(c) = 0\)

Then
\[
f \ ' (c) = \lim_{x\to c}{\frac{f(x)- f(c)}{x-c}} = \lim_{x\to c}{\frac{f(x)}{x-c}} = 0
\]

Also 
\[
g \ ' (c) = \lim_{x\to c}{\frac{|f(x)|}{x-c}}
\]
Then
\[
|g \ ' (c)| = \lim_{x\to c}{\left|\frac{|f(x)|}{x-c}\right|}=\lim_{x\to c}{\left|\frac{f(x)}{x-c}\right|} = |f \ ' (c)| = 0
\]

Hence \(g \ ' (c) = 0\)

\subsection*{BS - Page 167 Exercise 9}
Prove that if \(f:\mathbb{R} \rightarrow \mathbb{R}\) is an even function, and has a derivative at every point, then the derivative
\(f \ '\) is an odd function. Also prove that if \(g: \mathbb{R} \rightarrow \mathbb{R}\) is an odd differentiable function then
\(g \ '\) is an even function 
\\

\subsubsection*{(a)}
We have that \(f(-x) = f(x)\),

\[
(f(x))' = (f(-x))' 
\]
Then by the Chain Rule
\[
f \ ' (x) = f \ ' (-x)(-1) 
\]

\[
\Rightarrow f \ ' (-x) = - f \ ' (x)
\]
Thus \(f \ ' (x)\) is an odd function

\subsubsection*{(b)}
We have that \(g(-x) = -g(x)\),

\[
(g(-x))' = (-g(x))'
\]

Then by the Chain Rule
\[
g \ '(-x) (-1) = -g \ ' (x)
\]

\[
\Rightarrow g \ ' (-x) = g(x) 
\]
Thus \(g \ ' (x)\) is an even function

\subsection*{BS - Page 175 Exercise 11}
Give an example of a uniformly continuous function on \([0,1]\) that is differentiable on \((0,1)\), but whose derivative is 
not bounded on \((0,1)\)
\\

Consider \(f(x)= \sqrt{x}\), which I will prove is continuous on \([0,1]\)

Let \(\epsilon >0\)

Choose \(\delta = \epsilon^2 \)

Then for \(|x-c| < \delta\)
\begin{align*}
    |f(x)-f(c)|^2 &= |\sqrt{x} - \sqrt{c}|^2 \\
    &\leq |\sqrt{x} - \sqrt{c}||\sqrt{x} + \sqrt{c}| \\
    &< |x-c| \\
    &< \delta \\
    &< \epsilon^2
\end{align*}
Thus
\[
|f(x)- f(c)| < \epsilon 
\]

Since our choice of \(\delta\) was independent of \(c\), \(f(x)\) is uniformly continuous
\\

Now we look at \(f \ ' (x) = \frac{1}{2\sqrt{x}}\)

Let \(B >0\)

Assume \(f \ ' (x)\) is bounded

Then for all \(x \in (0,1)\), we have
\[
\frac{1}{2\sqrt{x}} \leq B 
\]
\[
\Rightarrow \sqrt{x} \geq \frac{1}{2B}
\]
But since \(\lim_{x\to0}{\sqrt{x}}=0\), we get that it can be arbitrarily small, thus 
for any given \(B>0\), \(\exists \ x\) s.t. the following inequality doesn't hold
\[
\sqrt{x} \geq \frac{1}{2B} > 0
\]
A contradiction, thus \(f \ ' (x)\) is unbounded

\end{document}